\section{Complete anatomy of the PSM voltage ellipse}
\label{sec:psm-ellipse}

\subsection{Shape matrix and exact center}

Squaring \cref{eq:psm-affine-map} gives
\begin{equation}
 \gls{sym:ip}^{\mathsf T}\gls{sym:hp}\gls{sym:ip}
 +2\vect b^{\mathsf T}\gls{sym:ip}+c\leq0,
 \end{equation}
where
\begin{align}
 \gls{sym:hp}&=\gls{sym:mp}^{\mathsf T}\gls{sym:mp}
 =\begin{bmatrix}
 R_s^2+\omega_e^2L_d^2&R_s\omega_e(L_d-L_q)\\
 R_s\omega_e(L_d-L_q)&R_s^2+\omega_e^2L_q^2
 \end{bmatrix},\label{eq:psm-Hp}\\
 \vect b&=\begin{bmatrix}\omega_e^2L_d\psi_{\mathrm{pm}}\\
 R_s\omega_e\psi_{\mathrm{pm}}\end{bmatrix},\qquad
 c=\omega_e^2\psi_{\mathrm{pm}}^2-U_{\max}^2.
\end{align}
The determinant
\begin{equation}
 \det\gls{sym:hp}=(R_s^2+\omega_e^2L_dL_q)^2
\end{equation}
is positive outside the joint ideal limit \(R_s=\omega_e=0\). Completing the
square therefore yields
\begin{equation}
 (\gls{sym:ip}-\gls{sym:ic})^{\mathsf T}\gls{sym:hp}
 (\gls{sym:ip}-\gls{sym:ic})=\gls{sym:umax}^2,
 \label{eq:psm-centered-ellipse}
\end{equation}
with
\begin{equation}
 \boxed{\gls{sym:ic}=\begin{bmatrix}
 -\dfrac{\psi_{\mathrm{pm}}\omega_e^2L_q}
 {R_s^2+\omega_e^2L_dL_q}\\[.7em]
 -\dfrac{\psi_{\mathrm{pm}}R_s\omega_e}
 {R_s^2+\omega_e^2L_dL_q}
 \end{bmatrix}.}
 \label{eq:psm-center}
\end{equation}
The alternative derivation is more physical: solve
\(\gls{sym:mp}\gls{sym:ic}+\vect r=\vect0\). The center is exactly the
stationary current that makes both stator-voltage components zero. It is not,
at finite speed and resistance, the familiar point
\((-\psi_{\mathrm{pm}}/L_d,0)\). That characteristic flux-cancellation point
is the high-speed limit and the exact ideal-resistance center.

\subsection{The center traces a fixed conic}

Substitution of \cref{eq:psm-center} gives the identity
\begin{equation}
 L_di_{\mathrm c,d}^2+L_qi_{\mathrm c,q}^2
 +\psi_{\mathrm{pm}}i_{\mathrm c,d}=0,
 \label{eq:center-locus-raw}
\end{equation}
or, after completing one scalar square,
\begin{equation}
 \boxed{L_d\left(i_{\mathrm c,d}+\frac{\psi_{\mathrm{pm}}}{2L_d}\right)^2
 +L_qi_{\mathrm c,q}^2=\frac{\psi_{\mathrm{pm}}^2}{4L_d}.}
 \label{eq:center-locus}
\end{equation}
The center starts at the origin when \(\omega_e=0\), travels along the lower
half for positive speed under the present sign convention, and approaches
\((-\psi_{\mathrm{pm}}/L_d,0)\) as \(\omega_e\to+\infty\). Negative speed
reflects the path across the \(d\)-axis.

Introduce
\begin{equation}
 \gls{sym:xi}=\frac{\gls{sym:omega}\sqrt{\gls{sym:ld}\gls{sym:lq}}}
 {\gls{sym:rs}}.
\end{equation}
For \(R_s>0\), the same center reads
\begin{equation}
 i_{\mathrm c,d}=-\frac{\psi_{\mathrm{pm}}}{L_d}
 \frac{\xi^2}{1+\xi^2},\qquad
 i_{\mathrm c,q}=-\frac{\psi_{\mathrm{pm}}}{\sqrt{L_dL_q}}
 \frac{\xi}{1+\xi^2}.
 \label{eq:center-dimensionless}
\end{equation}
Resistance changes how quickly physical speed traverses the locus; it does not
change the locus itself. If \(L_d=L_q\), \cref{eq:center-locus} is a circle
centered at \((-\psi_{\mathrm{pm}}/(2L_d),0)\).

\begin{figure}[tb]
\centering
\begin{tikzpicture}
\begin{axis}[axis main,width=.78\linewidth,height=6.2cm,axis equal image,
 xlabel={$i_{\mathrm c,d}/I_{\mathrm{ch}}$},ylabel={$i_{\mathrm c,q}/I_{\mathrm{ch}}$},
 xmin=-1.1,xmax=.12,ymin=-.68,ymax=.68,legend pos=north east]
 \addplot[mapped curve,domain=-8:8,samples=181]
 ({-x^2/(1+x^2)},{-1.22*x/(1+x^2)});
 \addplot[only marks,mark=*,MidnightBlue] coordinates {(0,0) (-1,0)};
 \addplot[coordinate vector] coordinates {(-.20,-.49) (-.63,-.58)};
 \addplot[coordinate vector] coordinates {(-.20,.49) (-.63,.58)};
 \node[annotation,anchor=west] at (axis cs:.02,0){$\omega_e=0$};
 \node[annotation,anchor=east] at (axis cs:-1.02,0){$|\omega_e|\to\infty$};
 \legend{center locus}
\end{axis}
\end{tikzpicture}
\caption{Normalized PSM center trajectory for $L_q/L_d=2/3$. Positive and
negative speed traverse opposite halves of the same conic. Resistance changes
the speed parameterization, not the curve.}
\label{fig:center-path}
\end{figure}


\subsection{Rotation and when an angle stops meaning anything}

For a quadratic \(Ax^2+Bxy+Cy^2\), principal directions satisfy
\(\tan(2\theta)=B/(A-C)\). In the present case the mixed coefficient is
\(B=2R_s\omega_e(L_d-L_q)\). It exists only when resistance, rotation, and
saliency are all present. If \(L_d\neq L_q\), cancellation of that common
factor gives
\begin{equation}
 \boxed{\tan(2\theta)=\frac{2R_s}{\omega_e(L_d+L_q)}.}
 \label{eq:psm-angle}
\end{equation}
The nonsalient case must be handled before this division: for \(L_d=L_q\),
\(\gls{sym:hp}=(R_s^2+\omega_e^2L_d^2)\matr I\), the boundary is a circle,
and every direction is a principal direction. Formula
\cref{eq:psm-angle} would otherwise assign a spurious orientation to a
degenerate \(0/0\) eigenvector problem.

At \(R_s=0\), the axes are the \(dq\) axes. At high speed the angle approaches
zero. At low but nonzero speed a formal angle may approach \(45^\circ\), yet
the ellipse simultaneously approaches the resistance circle; its eigenvalues
coalesce and the orientation becomes ill-conditioned and physically
irrelevant. The correct message is not that the low-speed circle is strongly
rotated, but that a nearly circular object has no robustly observable axis.

\begin{figure}[tb]
\centering
\begin{tikzpicture}
\begin{axis}[axis main,width=.82\linewidth,height=6.3cm,axis equal image,
 xlabel={$i_d$},ylabel={$i_q$},xmin=-3,xmax=2.2,ymin=-1.7,ymax=1.7]
 \addplot[voltage curve,BrickRed!45,domain=0:360,samples=121]
 ({-.25+1.38*cos(x)*cos(38)-.95*sin(x)*sin(38)},
  {1.38*cos(x)*sin(38)+.95*sin(x)*cos(38)});
 \addplot[voltage curve,BrickRed!72,domain=0:360,samples=121]
 ({-1.0+1.18*cos(x)*cos(16)-.63*sin(x)*sin(16)},
  {1.18*cos(x)*sin(16)+.63*sin(x)*cos(16)});
 \addplot[voltage curve,BrickRed,domain=0:360,samples=121]
 ({-1.55+.78*cos(x)},{.39*sin(x)});
 \node[annotation] at (axis cs:1.25,1.25){low speed:\\nearly circular};
 \node[annotation] at (axis cs:-.2,-1.25){medium};
 \node[annotation] at (axis cs:-2.15,.75){high speed:\\axes align};
\end{axis}
\end{tikzpicture}
\caption{Resistance and saliency generate the mixed term and therefore the
tilt. At low speed the formal axis angle is poorly meaningful because the
shape approaches a resistance circle; at high speed the axes approach $dq$.}
\label{fig:psm-tilt}
\end{figure}


\subsection{Characteristic current versus a copper-loss circle}

For a true PSM with stator copper loss only,
\begin{equation}
 \gls{sym:pcu}=\frac32\gls{sym:rs}(i_d^2+i_q^2),\qquad
 \gls{sym:ipmax}=\sqrt{\frac{2\gls{sym:pcumax}}{3\gls{sym:rs}}}.
 \label{eq:psm-loss-circle}
\end{equation}
Define \(\gls{sym:ich}=\gls{sym:psipm}/\gls{sym:ld}\). The point
\((-I_{\mathrm{ch}},0)\) may lie inside, on, or outside the admissible loss
circle. If a separate peak-current limit \(\gls{sym:imax}\) is also present,
the admissible radius is \(\min(I_P,I_{\max})\); the same three geometric
cases use that smaller radius:
\begin{itemize}
 \item \(I_{\mathrm{ch}}<I_P\): high-speed voltage ellipses shrink around an
       interior point. The ideal electromagnetic model permits arbitrarily
       high speed; finite-power operation can persist while torque decays
       approximately as \(1/|\omega_e|\).
 \item \(I_{\mathrm{ch}}=I_P\): the limiting center is on the thermal circle.
       This is the tangential characteristic-current case with vanishing
       current margin.
 \item \(I_{\mathrm{ch}}>I_P\): the limiting center lies outside. Continuity
       implies a finite speed at which the voltage ellipse and loss circle
       cease to intersect.
\end{itemize}
Mechanical speed, rotor stress, iron loss, inverter switching, demagnetization,
and thermal transients make the first statement an ideal-model limit, not a
promise of a realizable infinite-speed drive.

\begin{figure}[tb]
\centering
\begin{tikzpicture}[scale=.82]
\foreach \x/\c/\lab in {0/-1.15/inside,4.4/-1.55/on,8.8/-2.00/outside}{
 \begin{scope}[xshift=\x cm]
  \draw[->] (-2.35,0)--(.65,0) node[right]{$i_d$};
  \draw[->] (0,-1.45)--(0,1.55) node[above]{$i_q$};
  \draw[loss curve] (0,0) circle (1.55);
  \draw[voltage curve] (\c,0) ellipse (.44 and .72);
  \fill[BrickRed] (\c,0) circle (1.5pt);
  \node[below] at (-.75,-1.75){\lab};
 \end{scope}}
\end{tikzpicture}
\caption{Characteristic current relative to the PSM copper-loss circle.
Inside permits an ideal infinite-speed limiting center, on is tangential, and
outside implies a finite terminal speed when the shrinking voltage ellipse no
longer intersects the admissible disk.}
\label{fig:characteristic-cases}
\end{figure}


\paragraph{Connection to classical operation.}
The current/loss circle, voltage ellipse, and shifted torque hyperbola are the
standard geometric ingredients of \gls{mtpa}, \gls{fw}, and \gls{mtpv}
\cite{morimoto1990,kim1997}. Their intersections are generically quartic,
which explains the analytical PMSM literature. The additional EESM field
coordinate does not merely add another quartic variable; it restores global
homogeneity and introduces a torque-neutral direction that a 2D slice cannot
show.
